\documentclass[11pt]{amsart}
\usepackage[T1]{fontenc}
\usepackage{lmodern,microtype,amsmath,amssymb,amsthm,mathtools,enumitem}
\usepackage[margin=1.1in]{geometry}
\usepackage[hidelinks]{hyperref}
\usepackage{xurl}

\newtheorem{theorem}{Theorem}[section]
\newtheorem{lemma}[theorem]{Lemma}
\newtheorem{corollary}[theorem]{Corollary}
\newtheorem{proposition}[theorem]{Proposition}
\theoremstyle{definition}
\newtheorem{definition}[theorem]{Definition}
\theoremstyle{remark}
\newtheorem{remark}[theorem]{Remark}
\numberwithin{equation}{section}

\newcommand{\Q}{\mathbb Q}
\newcommand{\R}{\mathbb R}
\newcommand{\K}{\mathbb K}

\hypersetup{pdftitle={The relaxation time is not the value of a flow program},pdfauthor={Yitzchak Shmalo}}

\title{The relaxation time is not the value of a flow program}
\author{Yitzchak Shmalo}
\address{Einstein Institute of Mathematics, The Hebrew University of Jerusalem, Jerusalem, Israel}
\email{yitzchak.shmalo@gmail.com}

\subjclass[2020]{60J10, 90C05, 90C20, 05C21}
\keywords{relaxation time, spectral gap, multicommodity flow, linear programming, negative type metrics}

\begin{document}
\begin{abstract}
Aldous and Fill ask whether the relaxation time of a reversible Markov chain has a characterization as exactly an infimum over flows. It has none of the kind their flow bounds have. The optimal value of a linear or convex quadratic program in flow variables, with coefficients rational in the transition probabilities, lies in the field generated by the transition probabilities, while the relaxation time of the simple random walk on the path with six vertices is $3+\sqrt5$. The relaxation time is nevertheless an optimum of exactly the objective that is dual to minimum congestion: the dual of the congestion program maximises over metrics, and replacing metrics by squared Euclidean distances gives twice the relaxation time.
\end{abstract}
\maketitle

\section{Introduction}

Let $P=(p_{ij})$ be an irreducible transition matrix on a finite set $V$, reversible with respect to its stationary law $\pi$, and let $1=\lambda_1>\lambda_2\ge\dots\ge\lambda_n$ be its eigenvalues. The relaxation time is $\tau_2=1/(1-\lambda_2)$. Write $c_{ij}=\pi_ip_{ij}=\pi_jp_{ji}$ for the conductance of the edge $\{i,j\}$, $E$ for the set of edges with $c_{ij}>0$, and
\[
\mathcal E(f)=\sum_{\{i,j\}\in E}c_{ij}(f_i-f_j)^2,\qquad
\operatorname{Var}_\pi(f)=\tfrac12\sum_{a,b\in V}\pi_a\pi_b(f_a-f_b)^2 ,
\]
so that $\tau_2=\sup_f \operatorname{Var}_\pi(f)/\mathcal E(f)$ over nonconstant $f$ \cite[Ch.~3]{AF}.

Chapter~4 of Aldous and Fill \cite{AF} bounds $\tau_2$ above by quantities defined as minima over flows, in the line of Jerrum and Sinclair \cite{JS}, Diaconis and Stroock \cite{DS} and Sinclair \cite{S}: a family of flows routes a demand $\pi_a\pi_b$ from each $a$ to each $b$, and a congestion functional of the family, measured against the conductances, is minimised. Their Open Problem~3.44 asks: ``Does there exist a characterization of the relaxation time as exactly an inf over flows?'' The source adds that ``it would be more elegant to derive such inequalities from some exact characterization''.

Some restriction on the shape of the characterization is needed for the question to have content, since a constant functional of the flows equal to $\tau_2$ is an infimum over flows. The natural restriction is the one the chapter's bounds satisfy: a congestion functional or an energy, minimised over flows subject to linear constraints. Programs of this kind are linear or convex quadratic programs in the flow variables, with data rational in the entries of $P$. The answer for them is negative.

\begin{theorem}\label{thm:main}
Let $\Phi$ be a map that assigns to each reversible chain $P$ on a finite graph the optimal value of a linear or convex quadratic program, in finitely many real variables, whose objective and constraint coefficients are rational functions with rational coefficients of the entries of $P$, and which has a finite optimum. Then $\Phi(P)$ lies in the field $\Q(P)$ generated by the entries of $P$. For the simple random walk on the path with six vertices, $\Q(P)=\Q$ and $\tau_2=3+\sqrt5$. Hence $\Phi(P)\ne\tau_2(P)$ for that chain, and no such $\Phi$ equals $\tau_2$ on all reversible chains.
\end{theorem}

The stationary law is itself in $\Q(P)^V$, being the solution of a nonsingular linear system with coefficients in $\Q(P)$, so coefficients built from $\pi$ and the conductances $c_{ij}$ are allowed. The path is not special: the simple random walk on the path with five vertices has $\tau_2=2+\sqrt2$, on the five-cycle $\tau_2=1+\sqrt5/5$, and in a sample of $60$ random connected graphs on five to seven vertices with integer edge weights from $1$ to $4$, the relaxation time of the weighted walk is irrational in $59$ (Section~\ref{sec:numerics}). Rational values occur when the second eigenvalue is rational, as for $K_4$ ($\tau_2=3/4$) and the Petersen graph ($\tau_2=3/2$).

The objective that is dual to minimum congestion does yield $\tau_2$, if metrics are replaced by squared Euclidean distances. For a function $d$ on $V\times V$ put
\[
N(d)=\sum_{a,b\in V}\pi_a\pi_b\,d(a,b),\qquad C(d)=\sum_{\{i,j\}\in E}c_{ij}\,d(i,j).
\]
Let $\kappa$ be the minimum, over all families of flows routing $\pi_a\pi_b$ from $a$ to $b$ for every ordered pair $a\ne b$, of the largest ratio of the total load on an edge to its conductance.

\begin{proposition}\label{prop:cones}
\begin{enumerate}[label=\textup{(\roman*)}]
\item $\kappa=\max\{N(d): d \text{ a semimetric on } V,\ C(d)=1\}$.
\item $2\tau_2=\max\{N(d): d(a,b)=\|x_a-x_b\|^2 \text{ for some } x\colon V\to\R^k,\ k\ge1,\ C(d)=1\}$.
\end{enumerate}
\end{proposition}

Part (i) is the flow--metric duality for uniform multicommodity flow \cite{LR,LLR}; part (ii) is the variational formula for $\tau_2$ written for vector-valued functions. The functions in (ii) are the semimetrics of negative type \cite{Sch}, which form a cone that is not polyhedral. The congestion program and the relaxation time thus optimise the same objective under the same normalisation over two different cones, and the program over the metric cone is linear while the one over the negative type cone is semidefinite. On the path with six vertices the two values are $\kappa=5$ and $2\tau_2=6+2\sqrt5$.

Theorem~\ref{thm:main} leaves open whether some natural nonlinear functional of flows has optimum exactly $\tau_2$. The proof applies to any quantity that is irrational on a chain with rational transition probabilities, and in the other direction it says nothing about quantities such as the maximal mean hitting time, which solve linear systems and are rational on rational chains.

\section{Programs with data in a subfield}

\begin{lemma}\label{lem:field}
Let $\K\subseteq\R$ be a subfield. Let $Q$ be a symmetric positive semidefinite $m\times m$ matrix, $q\in\K^m$, $A\in\K^{r\times m}$, $b\in\K^r$, with $Q$ having entries in $\K$, and suppose that $g(x)=\frac12x^\top Qx+q^\top x$ is bounded below on the nonempty polyhedron $F=\{x\in\R^m: Ax\le b\}$. Then $g$ attains its infimum on $F$ at a point of $\K^m$, and the infimum lies in $\K$.
\end{lemma}

Equality constraints are included as pairs of inequalities, and $Q=0$ gives linear programs.

\begin{proof}
A quadratic function bounded below on a nonempty polyhedron attains its infimum there \cite{FW}; let $x^*$ be a minimiser and $I$ the set of rows of $A$ active at $x^*$. Since $g$ is convex and $F$ is polyhedral, the Karush--Kuhn--Tucker conditions hold at $x^*$ with some $y^*\ge0$, and conversely any $(x,y)$ satisfying them makes $x$ a minimiser. Consider the polyhedron
\[
T=\bigl\{(x,y)\in\R^m\times\R^I:\ Qx+q+A_I^\top y=0,\ A_Ix=b_I,\ y\ge0,\ Ax\le b\bigr\}.
\]
It is defined by linear equalities and inequalities with coefficients in $\K$ and contains $(x^*,y^*)$. A nonempty polyhedron has a minimal nonempty face, which is the solution set of a system of linear equations obtained by turning some of the defining inequalities into equalities. That system has coefficients in $\K$ and a real solution, so by Gaussian elimination over $\K$ it has a solution in $\K^{m+|I|}$. That point $(\bar x,\bar y)$ lies in $T$, so $\bar x$ satisfies the Karush--Kuhn--Tucker conditions and minimises $g$ on $F$, and $g(\bar x)\in\K$.
\end{proof}

\begin{proof}[Proof of Theorem~\ref{thm:main}]
Fix a chain $P$ and set $\K=\Q(P)\subseteq\R$. The program defining $\Phi(P)$ has data in $\K$ and a finite optimum, so $\Phi(P)\in\K$ by Lemma~\ref{lem:field}. For the simple random walk on the path $0,1,\dots,5$ every transition probability is $1/2$ or $1$, so $\K=\Q$. The eigenvalues of this chain are $\cos(\pi k/5)$, $k=0,\dots,5$, as one checks with the eigenvectors $f_j=\cos(\pi kj/5)$. Hence $\lambda_2=\cos(\pi/5)=(1+\sqrt5)/4$ and $\tau_2=4/(3-\sqrt5)=3+\sqrt5$, which is irrational.
\end{proof}

The same eigenvectors give $\lambda_2=\cos(\pi/4)$ and $\tau_2=2+\sqrt2$ on the path with five vertices; on the five-cycle $\lambda_2=\cos(2\pi/5)=(\sqrt5-1)/4$ and $\tau_2=1+\sqrt5/5$.

\begin{remark}
The two flow functionals that appear in the bounds are both covered. A congestion functional is a maximum of linear forms in the flows and is minimised by a linear program after introducing one more variable. An energy $\sum_e f_e^2/c_e$ of a unit flow is a convex quadratic function, and its minimum, the effective resistance, is rational on rational chains. On the path with six vertices the end-to-end energy is $50$ and $\kappa=5$.
\end{remark}

\section{Two cones}

\begin{proof}[Proof of Proposition~\ref{prop:cones}]
(i) The congestion program is the linear program of minimising $t$ over flows $f^{ab}\ge0$ on the directed edges, one for each ordered pair $a\ne b$, each sending $\pi_a\pi_b$ from $a$ to $b$, subject to $\sum_{a\ne b}\bigl(f^{ab}_{ij}+f^{ab}_{ji}\bigr)\le t\,c_{ij}$ for every edge. Its dual is to maximise $\sum_{a\ne b}\pi_a\pi_b\,\delta_\ell(a,b)$ over edge lengths $\ell\ge0$ with $\sum_E c_{ij}\ell_{ij}=1$, where $\delta_\ell$ is the shortest path distance for the lengths $\ell$ \cite{LR,LLR}. The function $\delta_\ell$ is a semimetric with $C(\delta_\ell)=1$. Conversely, a semimetric $d$ with $C(d)=1$ defines lengths $\ell_{ij}=d(i,j)$ with the same normalisation, and $\delta_\ell\ge d$ by the triangle inequality, so $N(\delta_\ell)\ge N(d)$. The two maxima coincide, and by duality they equal $\kappa$.

(ii) For $x\colon V\to\R^k$ with coordinates $x^{(1)},\dots,x^{(k)}$ and $d(a,b)=\|x_a-x_b\|^2$,
\[
N(d)=2\sum_{s=1}^k\operatorname{Var}_\pi(x^{(s)})\le 2\tau_2\sum_{s=1}^k\mathcal E(x^{(s)})=2\tau_2\,C(d),
\]
by the variational formula applied to each coordinate. Equality holds for $k=1$ and $x$ an eigenvector of $\lambda_2$, normalised so that $C(d)=1$.
\end{proof}

By Schoenberg's theorem \cite{Sch}, the functions in (ii) are exactly the semimetrics of negative type. The cone of negative type semimetrics on $V$ is not polyhedral when $|V|\ge3$, since it is linearly isomorphic to the cone of positive semidefinite $(|V|-1)\times(|V|-1)$ matrices. Squared Euclidean distances need not satisfy the triangle inequality, and metrics need not be of negative type (the graph metric of $K_{2,3}$ is not), so neither cone contains the other. Passing from metrics to negative type semimetrics is the step behind the semidefinite relaxation of sparsest cut of Arora, Rao and Vazirani \cite{ARV}. Kahale \cite{K} showed that the best bound on the spectral gap given by flows with weighted paths is the value of a semidefinite program and differs from the spectral gap by at most a polylogarithmic factor in the number of states. Proposition~\ref{prop:cones} places $\kappa$ and $2\tau_2$ as values of one objective on two different cones. It does not give an inequality between them, and none holds in general from this description alone.

\section{Computations}\label{sec:numerics}

The program \texttt{code/checks.py} in the repository computes $\lambda_2$ and $\tau_2$ exactly from the characteristic polynomial of the transition matrix for the simple random walks on $P_5$, $P_6$, $C_5$, $K_4$ and the Petersen graph, compares them with the closed forms above, and reads off irrationality from the minimal polynomial of $\tau_2$ ($x^2-6x+4$ for $P_6$). It computes $\kappa$ and the end-to-end energy exactly on the two paths ($15/4$ and $32$ on $P_5$, $5$ and $50$ on $P_6$), and solves the congestion program on $C_5$ as a linear program, with value $12/5$. It checks Proposition~\ref{prop:cones}(ii) by evaluating $N(d)/C(d)$ at an eigenvector of $\lambda_2$, which returns $2\tau_2$, and at $4{,}000$ random configurations in dimensions $1$ to $4$ for each chain, none of which exceeds it. Finally, for $60$ random connected graphs on five to seven vertices with integer edge weights between $1$ and $4$, it finds $\tau_2$ irrational in $59$ cases.

\section*{Acknowledgements and funding}
The research presented in this paper was supported by the European Research
Council (ERC) under the European Union's Horizon Europe research and
innovation programme (grant agreement No.~101041711), by the Simons
Foundation as part of the Collaboration on the Mathematical and Scientific
Foundations of Deep Learning, by Heights Labs, by Convex Nexus Capital, by
the Israel Science Foundation (grant number 2258/19), and by the Israel
Science Foundation (ISF Grant 4101/25). Computational resources were
provided in part by Google Cloud through its education credits program
(credit 530550103).

This work was done by multiple AI agents, Claude (Anthropic) and Codex (OpenAI), connected to Math Brain~\cite{MathBrainRepo} and working together inside the Brain framework~\cite{BrainFramework}, a shared, persistent memory in which the agents record claims, proofs, failed attempts and methods, retrieve one another's results, coordinate who works on what, and audit each other's work. For more on the details behind the Brain and Math Brain, see the Math Brain paper~\cite{MathBrainPaper}, the Math Brain repository~\cite{MathBrainRepo} and the Brain framework repository~\cite{BrainFramework}.

\section*{Declaration of competing interest}
The author holds equity in Heights Labs and in Convex Nexus Capital and
has a management role at Convex Nexus Capital. These entities are included
in the funding declaration above.

\section*{Declaration of generative AI and AI-assisted technologies}
The paper comes out of Math Brain, an ongoing project on the use of AI systems for research mathematics. In September 2026 Claude models (Anthropic), running as research agents of that project, located Open Problem~3.44 in the source, proposed the rationality argument for linear programs with the path witness, and computed the flow values. In November 2026 Claude Opus~5.5 (Anthropic), working in Claude Code, checked the argument, extended it to convex quadratic programs through Lemma~\ref{lem:field}, proved Proposition~\ref{prop:cones}, wrote the verification program and wrote the paper. The author directed the work and is responsible for the content of the paper and for decisions concerning its dissemination.

\section*{Data availability}
No data were used. The code is available at\\ \url{https://github.com/yspennstate/relaxation-time-flow-programs}.

The paper is part of the Math Brain project, whose public repository is \url{https://github.com/yspennstate/math-brain-public}~\cite{MathBrainRepo}.

\begin{thebibliography}{99}

\bibitem{AF}
D.~Aldous and J.~A.~Fill, \emph{Reversible Markov Chains and Random Walks on Graphs}, unfinished monograph, 2002, recompiled 2014, \url{https://www.stat.berkeley.edu/~aldous/RWG/book.html}.

\bibitem{ARV}
S.~Arora, S.~Rao and U.~Vazirani, Expander flows, geometric embeddings and graph partitioning, \emph{J. ACM} \textbf{56} (2009), no.~2, Art.~5, 1--37.
\bibitem{DS}
P.~Diaconis and D.~Stroock, Geometric bounds for eigenvalues of Markov chains, \emph{Ann. Appl. Probab.} \textbf{1} (1991), 36--61.

\bibitem{FW}
M.~Frank and P.~Wolfe, An algorithm for quadratic programming, \emph{Naval Res. Logist. Quart.} \textbf{3} (1956), 95--110.

\bibitem{JS}
M.~Jerrum and A.~Sinclair, Approximating the permanent, \emph{SIAM J. Comput.} \textbf{18} (1989), 1149--1178.
\bibitem{K}
N.~Kahale, A semidefinite bound for mixing rates of Markov chains, \emph{Random Structures Algorithms} \textbf{11} (1997), 299--313.
\bibitem{LR}
T.~Leighton and S.~Rao, Multicommodity max-flow min-cut theorems and their use in designing approximation algorithms, \emph{J. ACM} \textbf{46} (1999), 787--832.

\bibitem{LLR}
N.~Linial, E.~London and Y.~Rabinovich, The geometry of graphs and some of its algorithmic applications, \emph{Combinatorica} \textbf{15} (1995), 215--245.

\bibitem{Sch}
I.~J.~Schoenberg, Metric spaces and positive definite functions, \emph{Trans. Amer. Math. Soc.} \textbf{44} (1938), 522--536.

\bibitem{S}
A.~Sinclair, Improved bounds for mixing rates of Markov chains and multicommodity flow, \emph{Combin. Probab. Comput.} \textbf{1} (1992), 351--370.

\bibitem{code}
Y.~Shmalo, Code for ``The relaxation time is not the value of a flow program'', \url{https://github.com/yspennstate/relaxation-time-flow-programs}, 2026.

\bibitem{MathBrainRepo} Y.~Shmalo, \emph{Math Brain: public snapshot, code and tests}, 2026, \url{https://github.com/yspennstate/math-brain-public}.

\bibitem{MathBrainPaper} Y.~Shmalo, \emph{Math Brain: shared, auditable memory for multi-agent mathematical research}, October 2026, \url{https://github.com/yspennstate/math-brain-public/tree/main/papers/math-brain-2026-10-24}.

\bibitem{BrainFramework} Y.~Shmalo, \emph{Brain Framework}, public software repository, 2026, \url{https://github.com/yspennstate/brain-framework-public}.

\end{thebibliography}
\end{document}
